\documentclass[10pt,a4paper]{amsart}
\usepackage[margin=19mm]{geometry}
\usepackage{amsmath,amssymb,mathtools}
\usepackage[scr=rsfs,cal=euler]{mathalfa}
\usepackage{xfrac}
\usepackage[unicode,hidelinks,bookmarks=false]{hyperref}
\newcommand{\ie}{i.e.\ }
\newcommand{\cf}{cf.\ }
\newcommand{\resp}{respectively\ }
\newcommand{\Subsection}{\subsection}
\providecommand{\nobreakdash}{\nobreak\hbox{-}}
\setlength{\emergencystretch}{2em}
\multlinegap0pt
\usepackage{bm}
\usepackage{bbm}
\usepackage{dsfont}
\usepackage{xspace}
\usepackage{cancel}
\usepackage{mathtools}
\usepackage{cleveref}
\usepackage{amsthm}
\makeatletter
\@ifundefined{theorem}{\newtheorem{theorem}{Theorem}}{}
\makeatother
\DeclareMathOperator*{\argmax}{arg\,max}
\newcommand{\R}{\mathbb{R}}
\newcommand{\calC}{\mathcal{C}}
\newcommand{\calH}{\mathcal{H}}
\newcommand{\calM}{\mathcal{M}}
\newcommand{\calS}{\mathcal{S}}
\newcommand{\dd}{\,\mathrm{d}}
\newcommand{\pos}[1]{\left[#1\right]_+}
\title[Projective Maximum Entropy: Universality and Acceptance-Regi]{Projective Maximum Entropy: Universality and Acceptance-Region Calibration}
\author{Hideitsu Hino}
\date{Source: arXiv:2607.16547v1 (CC BY 4.0)}
\begin{document}
\maketitle

\noindent\emph{Source and licence.} Projective Maximum Entropy: Universality and Acceptance-Region Calibration. Version arXiv:2607.16547v1 (CC BY 4.0), distributed under the Creative Commons Attribution 4.0 International licence, as recorded in the arXiv licence metadata for this exact version and shown on the abstract page https://arxiv.org/abs/2607.16547v1. Full rights notice: licenses/arxiv-2607.16547-CC-BY-4.0.txt.\par\smallskip

\noindent\textit{Editorial scope.} Counted unit: the Theorem whose source label is \texttt{thm:qexp} in the article (source0.tex lines 415--426, source label \texttt{thm:qexp}), delivered together with the article's own complete original proof of it (source0.tex lines 428--443, one \texttt{proof} environment of 16 lines). No mathematics is written, repaired or re-derived: statement and proof are the article's own text line for line. Required context retained uncounted: eq:feasible-set (lines 192--204); thm:universality (lines 362--375). Every definition, display and label of those lines is retained unchanged. Omitted from this package and not redistributed: 1--191, 205--361, 376--414, 427--427, 444--1259. Nothing inside a retained excerpt is omitted or altered: every retained body is delivered exactly as the article's lines are written, so no omission receipt of any kind accompanies this package. Citations: the retained excerpts carry no citation command at all, so no bibliography entry of the article is transcribed and no reference is redistributed. Reference adaptation: none was needed and none was made; every reference inside the delivered unit resolves inside it, so no unresolved reference or citation is produced. Printed slips kept literal: no printed word, symbol, hypothesis or number of the counted statement or of its proof was altered. Layout and class: the article's own class is \texttt{sn-jnl} with packages including graphicx, multirow, amsmath, amssymb, amsfonts, amsthm, array, mathrsfs, appendix, xcolor, textcomp, manyfoot, booktabs, algorithm; the wrapper is the lane's pinned \texttt{amsart} editorial wrapper at 10pt on A4, which restates the theorem-like environments the retained text needs and adds the article's own macro definitions that the retained text needs (\texttt{\textbackslash R}, \texttt{\textbackslash argmax}, \texttt{\textbackslash calC}, \texttt{\textbackslash calH}, \texttt{\textbackslash calM}, \texttt{\textbackslash calS}, \texttt{\textbackslash dd}, \texttt{\textbackslash pos}), with extra wrapper packages (bm, bbm, dsfont, xspace, cancel, mathtools, cleveref). The delivered numbering is the unit's own. Assets: no retained span contains a figure environment, a graphics-inclusion command, a PDF-page inclusion, a table or a TikZ picture; no image file is copied into this package, the package declares no asset dependency and no image file is redistributed with it. Licence: version arXiv:2607.16547v1 of this article is distributed under the Creative Commons Attribution 4.0 International licence, as recorded in the arXiv licence metadata for this exact version and shown on the abstract page https://arxiv.org/abs/2607.16547v1. A full rights notice is delivered at licenses/arxiv-2607.16547-CC-BY-4.0.txt. Characters: every retained excerpt is pure ASCII, so the delivered engine, XeLaTeX with its default Latin Modern text font, reports no missing character.
\begin{equation}
\calC_{\tau,\gamma}
=
\left\{
 p\in\calM_{+,\gamma}:
 \int p\,\dd\mu_0=1,
 \quad
 \int \lvert T_j\rvert p\,\dd\mu_0<\infty,
 \quad
 \int Tp\,\dd\mu_0=\tau
\right\}.
\label{eq:feasible-set}
\end{equation}

\begin{theorem}[Projective maximum-entropy universality]\label{thm:universality}
Let $\gamma\in(-1,\infty)\setminus\{0\}$ and let $\calC_{\tau,\gamma}$ be defined by \cref{eq:feasible-set}. For every $\gamma$-admissible projective entropy,
\begin{equation}
 \argmax_{p\in\calC_{\tau,\gamma}}
 \calH_{\gamma,F}(p)
 =
 \argmax_{p\in\calC_{\tau,\gamma}}
 \calS_\gamma(p),
 \qquad
 \calS_\gamma(p)=\frac{1-\int p^{1+\gamma}\,\dd\mu_0}{\gamma}.
\label{eq:universality-argmax}
\end{equation}
If $\calC_{\tau,\gamma}$ is convex and an optimizer exists, the optimizer is unique.
\end{theorem}

\begin{theorem}[$q$-exponential representation]\label{thm:qexp}
Let $\gamma\in(-1,\infty)\setminus\{0\}$, and suppose that a maximum-entropy solution $p^*$ exists in $\calC_{\tau,\gamma}$ and satisfies the usual variational and Karush--Kuhn--Tucker regularity conditions. Then there exist $a\in\R$ and $\lambda\in\R^m$ such that
\begin{equation}
 p^*(x)=
 \begin{cases}
 \pos{a+\lambda^\top T(x)}^{1/\gamma},&\gamma>0,\\[1mm]
 \{a+\lambda^\top T(x)\}^{1/\gamma},&-1<\gamma<0,
 \end{cases}
\label{eq:qexp-general}
\end{equation}
where the bracket in the negative-power case must remain positive on the support of a finite interior solution. The constants are determined by normalization and the moment constraints.
\end{theorem}

\begin{proof}
By \cref{thm:universality}, it is sufficient to maximize $\calS_\gamma$. The Lagrangian is
\[
 \mathcal L(p)
 =
 \frac{1-\int p^{1+\gamma}\,\dd\mu_0}{\gamma}
 -\alpha\left(\int p\,\dd\mu_0-1\right)
 -\lambda^\top\left(\int Tp\,\dd\mu_0-\tau\right).
\]
On the region where $p>0$, the first variation gives
\[
 -\frac{1+\gamma}{\gamma}p(x)^\gamma
 -\alpha-\lambda^\top T(x)=0.
\]
After reparameterizing the multipliers, this becomes $p(x)^\gamma=a+\lambda^\top T(x)$. For $\gamma>0$, complementary slackness for the nonnegativity constraint produces the positive part in \cref{eq:qexp-general}. For $-1<\gamma<0$, the exponent is negative, so the affine bracket must be positive wherever the density is finite.
\end{proof}

\end{document}
